\section{Effective Cartier divisors on Noetherian schemes}
\label{section-Noetherian-effective-Cartier}

\noindent
In the locally Noetherian setting most of the discussion of
effective Cartier divisors and regular sections simplifies somewhat.

\begin{lemma}
\label{lemma-regular-section-associated-points}
Let $X$ be a locally Noetherian scheme. Let $\mathcal{L}$ be an invertible
$\mathcal{O}_X$-module. Let $s \in \Gamma(X, \mathcal{L})$. Then $s$
is a regular section if and only if $s$ does not vanish in the associated
points of $X$.
\end{lemma}

\begin{proof}
Omitted. Hint: reduce to the affine case and $\mathcal{L}$ trivial
and then use Lemma \ref{lemma-regular-section-structure-sheaf} and
Algebra, Lemma \ref{algebra-lemma-ass-zero-divisors}.
\end{proof}

\begin{lemma}
\label{lemma-effective-Cartier-in-points}
Let $X$ be a locally Noetherian scheme. Let $D \subset X$ be a closed subscheme
corresponding to the quasi-coherent ideal sheaf
$\mathcal{I} \subset \mathcal{O}_X$.
\begin{enumerate}
\item If for every $x \in D$ the ideal
$\mathcal{I}_x \subset \mathcal{O}_{X, x}$
can be generated by one element, then $D$ is locally principal.
\item If for every $x \in D$ the ideal
$\mathcal{I}_x \subset \mathcal{O}_{X, x}$
can be generated by a single nonzerodivisor, then $D$ is an
effective Cartier divisor.
\end{enumerate}
\end{lemma}

\begin{proof}
Let $\Spec(A)$ be an affine neighbourhood of a point $x \in D$.
Let $\mathfrak p \subset A$ be the prime corresponding to $x$.
Let $I \subset A$ be the ideal defining the trace of $D$ on
$\Spec(A)$. Since $A$ is Noetherian (as $X$ is locally Noetherian)
the ideal $I$ is generated by finitely many elements, say
$I = (f_1, \ldots, f_r)$. Under the assumption of (1) we have
$I_\mathfrak p = (f)$ for some $f \in A_\mathfrak p$.
Then $f_i = g_i f$ for some $g_i \in A_\mathfrak p$.
Write $g_i = a_i/h_i$ and $f = f'/h$ for some
$a_i, h_i, f', h \in A$, $h_i, h \not \in \mathfrak p$.
Then $I_{h_1 \ldots h_r h} \subset A_{h_1 \ldots h_r h}$ is
principal, because it is generated by $f'$. This proves (1).
For (2) we may assume $I = (f)$. The assumption implies
that the image of $f$ in $A_\mathfrak p$ is a nonzerodivisor.
Then $f$ is a nonzerodivisor on a neighbourhood of $x$ by
Algebra, Lemma \ref{algebra-lemma-regular-sequence-in-neighbourhood}.
This proves (2).
\end{proof}

\begin{lemma}
\label{lemma-effective-Cartier-codimension-1}
Let $X$ be a locally Noetherian scheme.
\begin{enumerate}
\item Let $D \subset X$ be a locally principal closed subscheme.
Let $\xi \in D$ be a generic point of an irreducible component of $D$.
Then $\dim(\mathcal{O}_{X, \xi}) \leq 1$.
\item Let $D \subset X$ be an effective Cartier divisor.
Let $\xi \in D$ be a generic point of an irreducible component of $D$.
Then $\dim(\mathcal{O}_{X, \xi}) = 1$.
\end{enumerate}
\end{lemma}

\begin{proof}
Proof of (1). By assumption we may assume $X = \Spec(A)$ and
$D = \Spec(A/(f))$ where $A$ is a Noetherian ring and $f \in A$.
Let $\xi$ correspond to the prime ideal $\mathfrak p \subset A$.
The assumption that $\xi$ is a generic point of an irreducible
component of $D$ signifies $\mathfrak p$ is minimal over $(f)$.
Thus $\dim(A_\mathfrak p) \leq 1$ by
Algebra, Lemma \ref{algebra-lemma-minimal-over-1}.

\medskip\noindent
Proof of (2). By part (1) we see that $\dim(\mathcal{O}_{X, \xi}) \leq 1$.
On the other hand, the local equation $f$ is a nonzerodivisor in
$A_\mathfrak p$ by Lemma \ref{lemma-characterize-effective-Cartier-divisor}
which implies the dimension is at least $1$ (because there must be a
prime in $A_\mathfrak p$ not containing $f$ by the elementary
Algebra, Lemma \ref{algebra-lemma-Zariski-topology}).
\end{proof}

\begin{lemma}
\label{lemma-integral-effective-Cartier-divisor-dvr}
Let $X$ be a Noetherian scheme. Let $D \subset X$ be an
integral closed subscheme which is also an
effective Cartier divisor. Then the local ring of $X$
at the generic point of $D$ is a discrete valuation ring.
\end{lemma}

\begin{proof}
By Lemma \ref{lemma-characterize-effective-Cartier-divisor}
we may assume $X = \Spec(A)$ and $D = \Spec(A/(f))$
where $A$ is a Noetherian ring and $f \in A$ is a nonzerodivisor.
The assumption that $D$ is integral signifies that $(f)$ is prime.
Hence the local ring of $X$ at the generic point is $A_{(f)}$
which is a Noetherian local ring whose maximal ideal is generated by
a nonzerodivisor. Thus it is a discrete valuation ring by
Algebra, Lemma \ref{algebra-lemma-characterize-dvr}.
\end{proof}

\begin{lemma}
\label{lemma-effective-Cartier-divisor-Sk}
Let $X$ be a locally Noetherian scheme. Let $D \subset X$ be an
effective Cartier divisor. If $X$ is $(S_k)$, then $D$ is $(S_{k - 1})$.
\end{lemma}

\begin{proof}
Let $x \in D$. Then $\mathcal{O}_{D, x} = \mathcal{O}_{X, x}/(f)$ where
$f \in \mathcal{O}_{X, x}$ is a nonzerodivisor. By assumption we have
$\text{depth}(\mathcal{O}_{X, x}) \geq \min(\dim(\mathcal{O}_{X, x}), k)$.
By Algebra, Lemma \ref{algebra-lemma-depth-drops-by-one} we have
$\text{depth}(\mathcal{O}_{D, x}) = \text{depth}(\mathcal{O}_{X, x}) - 1$
and by Algebra, Lemma \ref{algebra-lemma-one-equation}
$\dim(\mathcal{O}_{D, x}) = \dim(\mathcal{O}_{X, x}) - 1$.
It follows that
$\text{depth}(\mathcal{O}_{D, x}) \geq \min(\dim(\mathcal{O}_{D, x}), k - 1)$
as desired.
\end{proof}

\begin{lemma}
\label{lemma-normal-effective-Cartier-divisor-S1}
Let $X$ be a locally Noetherian normal scheme. Let $D \subset X$ be an
effective Cartier divisor. Then $D$ is $(S_1)$.
\end{lemma}

\begin{proof}
By Properties, Lemma \ref{properties-lemma-criterion-normal}
we see that $X$ is $(S_2)$. Thus we conclude by
Lemma \ref{lemma-effective-Cartier-divisor-Sk}.
\end{proof}

\begin{lemma}
\label{lemma-weil-divisor-is-cartier-UFD}
Let $X$ be a locally Noetherian scheme. Let $D \subset X$ be an integral
closed subscheme. Assume that
\begin{enumerate}
\item $D$ has codimension $1$ in $X$, and
\item $\mathcal{O}_{X, x}$ is a UFD for all $x \in D$.
\end{enumerate}
Then $D$ is an effective Cartier divisor.
\end{lemma}

\begin{proof}
Let $x \in D$ and set $A = \mathcal{O}_{X, x}$. Let $\mathfrak p \subset A$
correspond to the generic point of $D$. Then $A_\mathfrak p$ has dimension
$1$ by assumption (1). Thus $\mathfrak p$ is a prime ideal of height $1$.
Since $A$ is a UFD this implies that $\mathfrak p = (f)$ for some $f \in A$.
Of course $f$ is a nonzerodivisor and we conclude by
Lemma \ref{lemma-effective-Cartier-in-points}.
\end{proof}

\begin{lemma}
\label{lemma-codim-1-part}
Let $X$ be a Noetherian scheme. Let $Z \subset X$ be a closed subscheme.
Assume there exist
\begin{enumerate}
\item a collection of integral effective Cartier divisors
$D_i \subset X$, $i \in I$
\item a closed subset $Z' \subset X$ all of whose irreducible components
have codimension $\geq 2$ in $X$
(Topology, Definition \ref{topology-definition-codimension})
\end{enumerate}
such that $Z \subset Z' \cup \bigcup_{i \in I} D_i$ set-theoretically.
Then there exist integers $a_i \geq 0$ with $a_i = 0$ for almost all $i \in I$
such that the effective Cartier divisor
$$
D = \sum a_i D_i
$$
is contained in $Z$ and such that the inclusion morphism $D \to Z$
is an isomorphism away from codimension $2$ in $X$
(in the sense that there exists an open $U \subset Z$ such that
$D \cap U \to Z \cap U$ is an isomorphism and such that every
irreducible component of $Z \setminus U$ has codimension $\geq 2$ in $X$).
When $Z$ is nowhere dense in $X$ existence of the $D_i$, $i \in I$ and $Z'$
is guaranteed if $\mathcal{O}_{X, x}$ is a UFD for all $x \in Z$
or if $X$ is regular.
\end{lemma}

\begin{proof}
Let $\xi_i \in D_i$ be the generic point and let
$\mathcal{O}_i = \mathcal{O}_{X, \xi_i}$ be the local ring
which is a discrete valuation ring by
Lemma \ref{lemma-integral-effective-Cartier-divisor-dvr}.
Let $a_i \geq 0$ be the minimal valuation of an element of
$\mathcal{I}_{Z, \xi_i} \subset \mathcal{O}_i$.
Of course $a_i > 0$ only if $D_i$ is an irreducible component
of $Z$ and hence $a_i > 0$ only for a finite number of $i \in I$.
We claim that the effective Cartier divisor $D = \sum a_i D_i$ works.

\medskip\noindent
Namely, suppose that $x \in X$. Let $A = \mathcal{O}_{X, x}$.
Let $D_1, \ldots, D_n$ be the pairwise distinct divisors
$D_i$ such that $x \in D_i$ and $a_i > 0$.
For $1 \leq i \leq n$ let $f_i \in A$ be a local equation for $D_i$.
Then $f_i$ is a prime element of $A$ and $\mathcal{O}_i = A_{(f_i)}$. Let
$I = \mathcal{I}_{Z, x} \subset A$ be the stalk of the
ideal sheaf of $Z$. By our choice of $a_i$ we have
$I A_{(f_i)} = f_i^{a_i}A_{(f_i)}$. We claim that
$I \subset (\prod_{i = 1, \ldots, n} f_i^{a_i})$.

\medskip\noindent
Proof of the claim. The localization map
$\varphi : A/(f_i) \to A_{(f_i)}/f_iA_{(f_i)}$ is injective as
the prime ideal $(f_i)$ is the inverse image of the maximal ideal
$f_iA_{(f_i)}$. By induction on $n$ we deduce that
$\varphi_n : A/(f_i^n)\to A_{(f_i)}/f_i^nA_{(f_i)}$ is also injective.
Since $\varphi_{a_i}(I) = 0$, we have $I \subset (f_i^{a_i})$.
Thus, for any $x \in I$, we may write $x = f_1^{a_1}x_1$
for some $x_1 \in A$. Since $D_1, \ldots, D_n$ are pairwise
distinct, $f_i$ is a unit in $A_{(f_j)}$ for $i \not = j$.
Comparing $x$ and $x_1$ at $A_{(f_i)}$
for $n \geq i > 1$, we still have $x_1 \in (f_i^{a_i})$.
Repeating the previous process, we inductively write
$x_i = f_{i + 1}^{a_{i + 1}}x_{i + 1}$ for any $n > i \geq 1$.
In conclusion, $x \in (\prod_{i = 1, \ldots n} f_i^{a_i})$
for any $x \in I$ as desired.

\medskip\noindent
The claim shows that $\mathcal{I}_Z \subset \mathcal{I}_D$, i.e., that
$D \subset Z$. Moreover, we also see that $D$ and $Z$ agree at the $\xi_i$,
which proves that $D \to Z$ is an isomorphism away from codimension $2$ on $X$.

\medskip\noindent
To see the final statements we argue as follows. A regular local
ring is a UFD (More on Algebra, Lemma
\ref{more-algebra-lemma-regular-local-UFD}) hence it suffices
to argue in the UFD case. In that case, let
$D_i$ be the irreducible components of $Z$
which have codimension $1$ in $X$.
By Lemma \ref{lemma-weil-divisor-is-cartier-UFD} each $D_i$
is an effective Cartier divisor.
\end{proof}

\begin{lemma}
\label{lemma-codimension-1-is-effective-Cartier}
Let $Z \subset X$ be a closed subscheme of a Noetherian scheme. Assume
\begin{enumerate}
\item $Z$ has no embedded points,
\item every irreducible component of $Z$ has codimension $1$ in $X$,
\item every local ring $\mathcal{O}_{X, x}$, $x \in Z$ is
a UFD or $X$ is regular.
\end{enumerate}
Then $Z$ is an effective Cartier divisor.
\end{lemma}

\begin{proof}
Let $D = \sum a_i D_i$ be as in Lemma \ref{lemma-codim-1-part}
where $D_i \subset Z$ are the irreducible components of $Z$.
If $D \to Z$ is not an isomorphism, then $\mathcal{O}_Z \to \mathcal{O}_D$
has a nonzero kernel sitting in codimension $\geq 2$. This
would mean that $Z$ has embedded points, which is forbidden
by assumption (1). Hence $D \cong Z$ as desired.
\end{proof}

\begin{lemma}
\label{lemma-UFD-one-equation-CM}
Let $R$ be a Noetherian UFD. Let $I \subset R$ be an ideal
such that $R/I$ has no embedded primes and such that
every minimal prime over $I$ has height $1$.
Then $I = (f)$ for some $f \in R$.
\end{lemma}

\begin{proof}
By Lemma \ref{lemma-codimension-1-is-effective-Cartier}
the ideal sheaf $\tilde I$ is invertible on $\Spec(R)$.
By More on Algebra, Lemma \ref{more-algebra-lemma-UFD-Pic-trivial}
it is generated by a single element.
\end{proof}

\begin{lemma}
\label{lemma-effective-Cartier-divisor-is-a-sum}
Let $X$ be a Noetherian scheme. Let $D \subset X$ be an effective
Cartier divisor. Assume that there exist integral effective Cartier
divisors $D_i \subset X$ such that $D \subset \bigcup D_i$
set theoretically. Then $D = \sum a_i D_i$ for some $a_i \geq 0$.
The existence of the $D_i$ is guaranteed if $\mathcal{O}_{X, x}$
is a UFD for all $x \in D$ or if $X$ is regular.
\end{lemma}

\begin{proof}
Choose $a_i$ as in Lemma \ref{lemma-codim-1-part} and set $D' = \sum a_i D_i$.
Then $D' \to D$ is an inclusion of effective Cartier divisors which
is an isomorphism away from codimension $2$ on $X$. Pick $x \in X$.
Set $A = \mathcal{O}_{X, x}$ and let $f, f' \in A$ be the nonzerodivisor
generating the ideal of $D, D'$ in $A$. Then $f = gf'$ for some $g \in A$.
Moreover, for every prime $\mathfrak p$ of height $\leq 1$ of $A$ we see
that $g$ maps to a unit of $A_\mathfrak p$. This implies that $g$ is
a unit because the minimal primes over $(g)$ have height $1$
(Algebra, Lemma \ref{algebra-lemma-minimal-over-1}).
\end{proof}

\begin{lemma}
\label{lemma-quasi-projective-Noetherian-pic-effective-Cartier}
\begin{slogan}
On a projective scheme, every line bundle has a regular meromorphic section.
\end{slogan}
Let $X$ be a Noetherian scheme which has an ample invertible sheaf.
Then every invertible $\mathcal{O}_X$-module is isomorphic to
$$
\mathcal{O}_X(D - D') =
\mathcal{O}_X(D) \otimes_{\mathcal{O}_X} \mathcal{O}_X(D')^{\otimes -1}
$$
for some effective Cartier divisors $D, D'$ in $X$. Moreover, given a
finite subset $E \subset X$ we may choose $D, D'$ such that
$E \cap D = \emptyset$ and $E \cap D' = \emptyset$. If
$X$ is quasi-affine, then we may choose $D' = \emptyset$.
\end{lemma}

\begin{proof}
Let $x_1, \ldots, x_n$ be the associated points of $X$
(Lemma \ref{lemma-finite-ass}).

\medskip\noindent
If $X$ is quasi-affine and $\mathcal{N}$ is any invertible
$\mathcal{O}_X$-module, then we can pick a section $t$ of
$\mathcal{N}$ which does not vanish at any of the points
of $E \cup \{x_1, \ldots, x_n\}$, see Properties, Lemma
\ref{properties-lemma-quasi-affine-invertible-nonvanishing-section}.
Then $t$ is a regular section of $\mathcal{N}$ by
Lemma \ref{lemma-regular-section-associated-points}.
Hence $\mathcal{N} \cong \mathcal{O}_X(D)$ where
$D = Z(t)$ is the effective Cartier divisor corresponding to $t$, see
Lemma \ref{lemma-characterize-OD}. Since $E \cap D = \emptyset$
by construction we are done in this case.

\medskip\noindent
Returning to the general case, let $\mathcal{L}$ be an ample invertible sheaf
on $X$. There exists an $m > 0$ and a section
$s \in \Gamma(X, \mathcal{L}^{\otimes m})$ such that $X_s$
is affine and such that $E \cup \{x_1, \ldots, x_n\} \subset X_s$
(Properties, Lemma \ref{properties-lemma-ample-finite-set-in-principal-affine}).

\medskip\noindent
Let $\mathcal{N}$ be an arbitrary invertible $\mathcal{O}_X$-module.
By the quasi-affine case, we can find a section
$t \in \mathcal{N}(X_s)$ which does not vanish at any point
of $E \cup \{x_1, \ldots, x_n\}$.
By Properties, Lemma \ref{properties-lemma-invert-s-sections}
we see that for some $e \geq 0$ the section $s^e|_{X_s} t$ extends to
a global section $\tau$ of $\mathcal{L}^{\otimes me} \otimes \mathcal{N}$.
Thus both $\mathcal{L}^{\otimes me} \otimes \mathcal{N}$ and
$\mathcal{L}^{\otimes me}$ are invertible sheaves which have global sections
which do not vanish at any point of $E \cup \{x_1, \ldots, x_n\}$.
Thus these are regular sections by
Lemma \ref{lemma-regular-section-associated-points}.
Hence $\mathcal{L}^{\otimes me} \otimes \mathcal{N} \cong \mathcal{O}_X(D)$
and $\mathcal{L}^{\otimes me} \cong \mathcal{O}_X(D')$ for some
effective Cartier divisors $D$ and $D'$, see Lemma \ref{lemma-characterize-OD}.
By construction $E \cap D = \emptyset$ and $E \cap D' = \emptyset$
and the proof is complete.
\end{proof}

\begin{lemma}
\label{lemma-wedge-product-ses}
Let $X$ be an integral regular scheme of dimension $2$.
Let $i : D \to X$ be the immersion of an effective Cartier divisor.
Let $\mathcal{F} \to \mathcal{F}' \to i_*\mathcal{G} \to 0$
be an exact sequence of coherent $\mathcal{O}_X$-modules.
Assume
\begin{enumerate}
\item $\mathcal{F}, \mathcal{F}'$ are locally free of rank $r$ on a nonempty
open of $X$,
\item $D$ is an integral scheme,
\item $\mathcal{G}$ is a finite locally free $\mathcal{O}_D$-module
of rank $s$.
\end{enumerate}
Then $\mathcal{L} = (\wedge^r\mathcal{F})^{**}$ and
$\mathcal{L}' = (\wedge^r \mathcal{F}')^{**}$
are invertible $\mathcal{O}_X$-modules and
$\mathcal{L}' \cong \mathcal{L}(k D)$ for some
$k \in \{0, \ldots, \min(s, r)\}$.
\end{lemma}

\begin{proof}
The first statement follows from Lemma \ref{lemma-reflexive-over-regular-dim-2}
as assumption (1) implies that $\mathcal{L}$ and $\mathcal{L}'$
have rank $1$. Taking $\wedge^r$ and double duals are functors, hence
we obtain a canonical map $\sigma : \mathcal{L} \to \mathcal{L}'$
which is an isomorphism over the nonempty open of (1), hence
nonzero. To finish the proof, it suffices to see that
$\sigma$ viewed as a global section of
$\mathcal{L}' \otimes \mathcal{L}^{\otimes -1}$ does not
vanish at any codimension point of $X$, except at the generic
point of $D$ and there with vanishing order at most $\min(s, r)$.

\medskip\noindent
Translated into algebra, we arrive at the following problem:
Let $(A, \mathfrak m, \kappa)$ be a discrete valuation ring
with fraction field $K$. Let $M \to M' \to N \to 0$ be an exact sequence
of finite $A$-modules with $\dim_K(M \otimes K) = \dim_K(M' \otimes K) = r$
and with $N \cong \kappa^{\oplus s}$. Show that the induced map
$L = \wedge^r(M)^{**} \to L' = \wedge^r(M')^{**}$ vanishes to
order at most $\min(s, r)$. We will use the structure theorem for
modules over $A$, see
More on Algebra, Lemma
\ref{more-algebra-lemma-generalized-valuation-ring-modules} or
\ref{more-algebra-lemma-modules-PID}.
Dividing out a finite $A$-module by a torsion submodule does not
change the double dual.
Thus we may replace $M$ by $M/M_{tors}$ and $M'$ by
$M'/\Im(M_{tors} \to M')$ and assume that $M$ is torsion free.
Then $M \to M'$ is injective and $M'_{tors} \to N$ is injective.
Hence we may replace $M'$ by $M'/M'_{tors}$ and $N$ by $N/M'_{tors}$.
Thus we reduce to the case where $M$ and $M'$ are free of rank $r$
and $N \cong \kappa^{\oplus s}$. In this case $\sigma$
is the determinant of $M \to M'$ and vanishes to order $s$
for example by Algebra, Lemma \ref{algebra-lemma-order-vanishing-determinant}.
\end{proof}









